\chapter{Mastery and the Retrieval Loop}

\otbookepigraph{Measure capability, not study time.}{PLDS Principle}

\section{The Reconstruction Test}

After reading Spong's chapter on \otindex{Velocity Kinematics}, close the book immediately.

\begin{practice}[Reconstruction Challenge]
    Starting only from the forward kinematics \(x = f(q)\), derive the relationship for tip velocity \(\dot{x}\).
    \textit{Hint: Use the chain rule to reveal the \otindex{Jacobian} \(J\).}
\end{practice}

\section{Stress-Testing the Robot}

Robotics is a ``Mixed Project Domain''. A model that works for a 2-link planar arm must survive a \otindex{Stress-test}.

\begin{designrule}[Breaking the Model]
    Take your derived Jacobian and set the robot into a \otindex{Singularity} configuration.
    \begin{itemize}
        \item What happens to the inverse \(J^{-1}\)?
        \item Physically, why can't the robot move in a specific direction?
    \end{itemize}
\end{designrule}

\section{Closing the Loop}

Use the \otindex{3-Day-On / 3-Day-Off} cycle for the control laws in Spong Chapter 7:
\begin{itemize}
    \item \textbf{On Block:} Build the model of PD Control with Gravity Compensation.
    \item \textbf{Off Block:} Step away and work on your \otindex{LaTeX-lab} templates.
    \item \textbf{Return Cold:} Attempt to draw the block diagram of the control system from memory.
\end{itemize}

\begin{summarybox}
    You have mastered a chapter in Spong when you can:
    \begin{enumerate}
        \item Explain the necessity of the main terms.
        \item Rebuild the core derivation without looking.
        \item Predict where the model breaks (singularities, limits).
    \end{enumerate}
\end{summarybox}
